









\chapter{Apparatus and Sample}

Our experimental system can be divided into two parts: apparatus and sample. The apparatus includes typical NMR devices and detectors. 
To serve the dark matter measurement, we adopt some special designs of the apparatus as well. Moreover, we select NMR samples for calibrating the system and performing the measurement. 

\section{Overview of NMR apparatus}

Typical NMR apparatus can be divided into three categories by their functions: signal generation, reception and initial processing. The first category includes the magnet and excitation coil which generate precession of nuclear spins. The second one is essentially the detector to pick up spin-precession signal. NMR signal is likely to be weak and noisy, hence the last category of the apparatus is the devices for initial signal processing, such as amplifying. 

According to the procedure described in Fig.\,\ref{fig:Procedure}, we set NMR resonant frequency to different values in order to scan a large frequency band. Since the resonant frequency is proportional to the magnitude of leading field, magnet should be able to be ramped to different magnitudes, for which reason we rule out permanent magnet. 
To perform precise and accurate measurements, we prefer low-noise apparatus and environment. The field generated by normal-conducting coil can be set arbitrarily. However, inefficient cooling limits the current in the coil. And external current supply introduces noise into the detector. In contrast, superconducting magnet can be ramped to over \SI{10}{\tesla}. Once magnetic field is set, the field can be isolated from external electric noise. Considering these advantages, a superconducting magnet with maximal field of \SI{0.1}{\tesla} was chosen for the CASPEr-gradient low-field experiments.
% Shim coils were can manipulate the homogeneity of the leading field. 
On top of this, an excitation coil was installed to realize pulsed-NMR diagnostic. 
Besides, a \SI{14.1}{\tesla} superconducting magnet was ordered for CASPEr-gradient high-field phase.  

As to the detector of experiment, the combination of gradiometer and direct-current superconducting quantum interference device (dc SQUID) was chosen for the CASPEr-gradient low-field phase. For high-field phase, a pick-up coil and amplifier function as the detector. 
%This thesis focuses on the low-field experiments, so we do not discuss the details of the other phase here. 

To process the output from the detector, a lock-in amplifier (LIA) with built-in low-pass filters is adopted for noise cancellation and signal amplification. It is also able to convert analog signal to digital form for terminals like computer to read. 

\section{Cryostat and superconducting magnet}

The cryostat (see Fig.\,\ref{fig:Cryostat_and_SCdevices}(a)) is placed inside a Faraday room, shielded from the external electric field.  It consists of the outer shell, reservoirs for liquid nitrogen (L\ch{N2}) and liquid helium (L\ch{He}), and vacuum insulation layers in between. L\ch{N2} is meant to cool down the cryostat and reduce the evaporation of L\ch{He}. The \ch{Nb} shield and superconducting devices inside can be submerged underneath L\ch{He} bath. 
Once filled with L\ch{N2} and L\ch{He}, the cryostat can maintain a superconducting environment for one week, and \ch{Nb} shield can repel external magnetic field so as to insulate the devices inside it from external magnetic noise. 

\begin{figure}[htbp]
    \includegraphics[width = \textwidth]{figures/02_Apparatus/Cryostat_and_SCdevices_20220903.png}
    \caption{(a) Cryostat for maintaining superconducting temperature. L\ch{He} bath offers \SI{4.2}{\kelvin} environment for superconducting apparatus. 
    (b) Superconducting apparatus. Besides main coil, shim and excitation coils, SQUIDs with \ch{Nb} capsule, \ch{Nb} wires for connections, gradiometer made of \ch{Nb} wires require \SI{9}{\kelvin} for operation. }
    \label{fig:Cryostat_and_SCdevices}
\end{figure}


The magnet system consists of main and shim coils. The main magnet offers \SI{0}{} to \SI{0.1}{\tesla} leading field, while shim coils can adjust the homogeneity of it with 8 coils, which are so-called shim channels with names of Z1, Z2, X, Y, XY, XZ, YZ and X2-Y2. The channel name indicates the geometry of the magnetic field it generates. A brief summary of their properties is shown in Tab.\,\ref{tab:shimproperity}. 

\begin{table}[ht]
\centering
\caption{Properties of shim coils. The order is the exponent of the $\mathrm{B}_z$ field equation. The interaction type implies the sequence in shimming. In general, smaller number means earlier in the shimming sequence. This table is reproduced from \cite{Conover1999TheSO}. }
\begin{tabular}{cccc}
	\hline
	\multirow{2}{*}{Channel name} & Equation for $\mathrm{B}_z$ field generated & \multirow{2}{*}{Order} & \multirow{2}{*}{Interaction Type} \\
	 & (not to exact magnitude) & & \\
	\midrule
	Z0 & $1$ & 0 & 0 \\
	% \midrule
	Z1 & $z$ & 1 & 1 \\
	% \midrule
	Z2 & $2z^2 - (x^2 + y^2)$ & 2 & 1 \\
	% \midrule
	X & $x$ & 1 & 0 \\
	% \midrule
	Y & $y$ & 1 & 0 \\
	% \midrule
	ZX & $zx$ & 2 & 2 \\
	% \midrule
	ZY & $zy$ & 2 & 2 \\
	% \midrule
	XY & $xy$ & 2 & 1 \\
	% \midrule
	X2-Y2 & $x^2-y^2$ & 2 & 1 \\
	\bottomrule
\end{tabular}
% \footnotesize{} %  \textit{The Shimming of High Resolution NMR Magnets}
\label{tab:shimproperity}
\end{table}


\section{Excitation generation}

To perform NMR schemes we mentioned in \S\,\ref{sec:intro:NMR}, we put a excitation coil next to the sample (see Fig.\,\ref{fig:Cryostat_and_SCdevices}(b)). This coil is made of \ch{Nb} wires to avoid Johnson–Nyquist noise, which is generated by thermal agitation of charge carriers inside normal-conducting material like copper. Excitation coil comes with $\mathrm{B_x}$, $\mathrm{B_y}$, $\mathrm{B_z}$ and $\mathrm{d B_z/dz}$ channels. The first three channels can produce homogeneous field along x, y and z-axis respectively, while the last one can produce a linear gradient along z-axis. 

A function generator along with a radio-frequency (rf) amplifier generate the pulse for excitation. For transmitting the pulse, a tuning circuit made of capacitors and inductors is designed for impedance matching between function generator and excitation coil at different frequencies. 


\section{Signal reception and processing}\label{sec:Signal_receptio_and_processing}

When magnetic flux is generated by the sample in the gradiometer (a deliberately-designed pick-up coil), the flux couples to superconducting loops of SQUIDs, each of which has two Josephson junctions, and gets converted to voltage signal. The voltage output of SQUID goes into LIA. The analog frequency mixers and RC filters for the demodulation in LIA can amplifier signals from interesting frequency band set by the user. The noises outside the band can be filtered out. 
After initial processing in LIA, the signal can be read by the computer for further processing and analyses. 

In this section we discuss the layout and geometry of the gradiometer and sample holder. Furthermore, we present computation on the coupling between gradiometer and SQUID concerning the detector response. We do not discuss the technical details of LIA here. 

\subsection{Gradiometer and sample holder}\label{sec:Signal_receptio_and_processing:Gradiometer_and_sample_holder}

Sample of large volume is not preferred by regular NMR experiments due to inhomogenous broadening which obstacles acquiring high-resolution spectrum. However, when we are expecting coupling between axion gradient and our sample, large sample is desired. The broadening enables the magnetometer to probe axion in a larger frequency band compared to that of smaller sample. Detailed discussion on the trade between sensitivity and axion-probing frequency band is presented in \S\,\ref{sec:Scanning_strategy}. 
Due to the geometry of the cryostat and magnet, the coil and sample must be fitted into a cylindrical region. Therefore, the gradiometer and sample holder are both designed cylinder-like (see Fig.\,\ref{fig:Gradiometer_and_Sample}). 

Three saddle coils made of pure \ch{Nb} wires are connected in a row to function as the pick-up coil (check Fig.\,\ref{fig:Gradiometer_and_Sample}(b) for the schematic). 
Such loops are often called gradiometer since they are widely used in measuring spatial gradient of magnetic flux. 
The ratio of coil normal areas of are approximately -1:2:-1. 
Because of such design, sample at the center of gradiometer can generate flux in the central saddle coil, while those in side coils are smaller considering the larger distances between sample and side coils. Therefore the sample signal, or so-called near-field signal, do not cancel out itself. 
In contrast, consider the response in gradiometer generated by a noise source far away. The noise is effectively plane wave for the gradiometer. Since the sum of normal area is zero, the flux in side coils cancels the one in central coil. Hence the gradiometer is sensitive to near-field NMR signal and insensitive to far-field noise. 
\begin{figure}[htbp]
	\centering
	\includegraphics[width=0.8\textwidth]{figures/02_Apparatus/Gradiometer_and_Sample_20220903.png}
    \caption{(a) Simplified schematic of gradiometer and sample. Details of gradiometer loop connection is shown in (b). }
    \label{fig:Gradiometer_and_Sample}
\end{figure}

\subsection{SQUID detector}\label{sec:Signal_receptio_and_processing:SQUID}

Unlike a coil which transforms the time derivative of flux into voltage, SQUID makes use of flux quantization\cite{LondonSuperfluids, DeaverPhysRevLett.7.43, DollPhysRevLett.7.51} and Josephson tunneling \cite{GiaeverPhysRevLett.5.147,GiaeverPhysRevLett.5.464,JOSEPHSON1962251}, converting the flux coupled to the SQUID loop into current or voltage signal. The SQUID circuit is installed in a \ch{Nb} capsule for shielding and protection. 
\begin{figure}[t]
	\centering
	\includegraphics[width=\textwidth]{figures/02_Apparatus/SQUID_Schematic_VPhi_FLL.png}
    \caption{(a) $V-\Phi$ characteristic of dc SQUID. W stands for the working point of SQUID, where the $\delta V$ has a linear dependence on $\delta \Phi$ approximately .
    (b) Direct-coupled flux-lock loop (FLL) circuit of SQUID with built-in pre-amplifier and integrator. 
    Figure is taken from \cite{squidhandbookintro}.}
    \label{fig:SQUID_schematic}
\end{figure}

One possible scheme for flux-voltage conversion is illustrated in Fig.\,\ref{fig:SQUID_schematic}(a). 
We denote the flux in the gradiometer as $\Phi_{grad}$. The gradiometer is connected to input coil inside SQUID capsule. Here we name it pick-up coil loop to distinguish it from SQUID loop which has two Josephson junctions. 
The total inductance of the pick-up coil loop $L$ is the sum of gradiometer inductance $L_{grad}$ and SQUID input inductance $L_{in}$:
\begin{align}
    L = L_{grad} + L_{in}\,,\label{eq:L_Lgrad_Lin}
\end{align}
while the current $I$ in the pick-up coil loop is
\begin{align}
	I = \dfrac{\Phi_{grad}}{L}\,.\label{eq:I_Phigrad_L}
\end{align}
Such current generates flux $\Phi_{in}$ in SQUID loop, which is characterized by 
\begin{align}
	\Phi_{in} = I \mathrm{g}_{in}\,.\label{eq:Phiin_I_Min}
\end{align}
Here $1/\mathrm{g}_{in}$ is the input coupling strength. We can obtain the relationship between $\Phi_{in}$ and $\Phi_{grad}$ from Eqs.\,\eqref{eq:L_Lgrad_Lin}, \eqref{eq:I_Phigrad_L} and \eqref{eq:Phiin_I_Min}:
\begin{align}
	\Phi_{in} = \dfrac{\mathrm{g}_{in}}{L_{grad} + L_{in}}\Phi_{grad}\,.\label{eq:Phiin_Phigrad}
\end{align}

Flux quantization demands the flux in the SQUID loop $\Phi_{SQUID}$ to be integer number of flux quanta $\Phi_0 = h/2e$, where $h$ is Planck constant and $e$ is the electron charge. Therefore, feedback current $I_f$ and corresponding flux $-\Phi_{in}$ are generated in the loop to cancel $\Phi_{in}$ when at the working point illustrated in Fig.\,\ref{fig:SQUID_schematic}(a). In this way the flux quantization condition is satisfied. 
Since we do not have to determine the absolute orientation of the flux, we omit the minus symbol for brevity and have:
\begin{align}
	I_f = \dfrac{\Phi_{in}}{\mathrm{g}_{f}} \,, \label{eq:If_Phiin_Mf}
\end{align}
where $1/\mathrm{g}_{f}$ is the feedback sensitivity. Such a variance in current can be read out as voltage signal $V_f$ with a feedback resistance $R_f$:
\begin{align}
	V_f = I_f R_f\,, \label{eq:Vf_If_Rf}
\end{align}

 With Eqs.\,\eqref{eq:L_Lgrad_Lin}, \eqref{eq:I_Phigrad_L}, \eqref{eq:Phiin_I_Min}, \eqref{eq:If_Phiin_Mf}, and \eqref{eq:Vf_If_Rf}, we obtain the feedback voltage signal as
\begin{align}
	V_f =\dfrac{ R_f}{ \mathrm{g}_{f} }  \Phi_{in} = \dfrac{ R_f}{\mathrm{g}_{f} } \dfrac{\mathrm{g}_{in}}{L_{grad} + L_{in}}\Phi_{grad} \,.\label{eq:Vf_Phigrad}
\end{align}

By far we have acquired the equations for $\Phi_{in}-\Phi_{grad}$ and $V_f - \Phi_{in}$ characteristics, which provides essential information in describing signal signature. 
% \sisetup{exponent-mode = scientific}
\begin{table}[htb]
    \centering
    \caption{Properties of SQUID sensors in use. Note that $R_f$ can be set to different values. $L_{grad}$ is computed to be \SI{553}{\nano\henry}. }
    \begin{tabular}{cccccccc}
    	\hline
    	\multirow{2}{*}{Sensor ID} & $L_{in}$ & $\mathrm{g}_{in}$ & $\mathrm{g}_{f}$ & $R_f$ &  $R_f/\mathrm{g}_{f}$ & \multirow{2}{*}{$\dfrac{\mathrm{g}_{in}}{L_{grad} + L_{in}}$} \\
    	 & (\SI{}{\nano \henry}) & ($\Phi_{0}$/\SI{}{\mA}) & ($\Phi_{0}$/\SI{}{\mA}) & (\SI{}{\kohm}) & (\SI{}{\volt}/$\Phi_{0}$) & \\
    	\midrule
    % 	C6L1W & \SI{400}{} & 1/31.54 & 1/44.12  \\
    	C6L1W & \SI{400}{} & 31.71 & 22.67 & 10 & 0.441 & \SI{6.88e-5}{}\\
    	\midrule
    % 	C73L1 & \SI{400}{} & 1/30.66 & 1/44.20 \\
    	C73L1 & \SI{400}{} & 32.62 & 22.62 & 10 & 0.442 & \SI{7.14e-5}{}\\
    	\bottomrule
    \end{tabular}
    \label{tab:SQUIDsensors}
\end{table}
% \sisetup{exponent-mode = input}


\section{Experimental control}

NMR experiments contain complicated sequences for pulsing and signal reception. Besides, the timing of steps are vital at the scale of microsecond. To appropriately coordinate between devices made by various parties, instantly analyze data to monitor experiment progress and communicate feedback in the network of software and hardware, a centralized control based on Python was created. It is named ``CASPEr Code'' and internally available via GitLab. The basic functions of it include:
\begin{enumerate}
    \item Shim control
    \begin{enumerate}
        \item Remote switching the shim power on and off. 
        \item Remotely connecting or cutting off communication between shim power supply and coils. 
        \item Setting the shim currents to desired values. The heaters can also be turned on and off for setting the currents. 
    \end{enumerate}
    \item Pulse control, logging and data analysis
    \begin{enumerate}
        \item Command for injecting pulse. 
        \item Recording the timing of rf pulses by oscilloscope in HDF5 file. 
        \item Analysis on function generator data file. 
    \end{enumerate}
    \item LIA control and data processing
    \begin{enumerate}
        \item Recording SQUID data and creating HDF5 file for data storage. 
        \item Analysis on HDF5 file data. 
    \end{enumerate}
\end{enumerate}

The combination of basic functions can realized more complicated schemes. We list a few schemes that has proven useful in practice:
\begin{enumerate}
    \item Pulsed-NMR. 
    \item Shimming to manipulate NMR linewidth with feedback from pulsed-NMR diagnostic. 
    \item Axion wind measurement done with LIA recording. 
    \item Magnetic field and sample temperature monitoring done with repetition of pulsed-NMR diagnostic. 
\end{enumerate}
More schemes can be realized by assembling basic functions. 

In addition to the hardware and software control, ``CASPEr Code'' have features for NMR kinetic simulation, stochastic axion wind simulation, apparatus properties computation and etc. It also stores important files online and synchronizes over multiple terminals. 

\section{Sample}

Samples need to satisfy the requirements from different stages of experiment. For the axion measurement phase, we desire a sample that could generate strongest signal from gradient coupling. However, we may need to choose a different sample for earlier stage to test the apparatus and calibrate the relationship between external excitation and signal response. Therefore we divide our discussion into "research sample" and "calibration sample". 

\subsection{Research sample}

Spin density is defined as the density of nuclear spins. Since we expect axion to interact with nuclear spins in a sample holder of finite volume, spin density determines signal amplitude. In this sense, condensed-state samples have advantages over gas sample. 

In addition to the phase and molecular properties, polarization also plays an important role. 
%From Tab.\,\ref{tab:samplesproperity} we can find different choices of liquid chemical have maximal difference of two orders of magnitude in terms of spin density. However, a
As discussed in \S\,\ref{sec:intro:NMR:Zeeman_Pol}, hyper-polarization can boost the effective spin density by 6 orders of magnitude. Therefore, $\ch{^{129}Xe}$ with spin 1/2 was chosen as the research sample in that there has been mature technology to hyper-polarize it through spin exchange optically pumping (SEOP). 

In this thesis, we do not provide further details on SEOP and hyper-polarization. Instead, we demonstrate the whole process of axion measurement with thermally-polarized sample. 

\subsection{Calibration sample}

We can list requirements for the NMR calibration sample (ranked by priority):
\begin{itemize}
    \item The test done with calibration sample should be helpful for the axion measurement. Depending on the sample properties, different schemes of NMR experiment are adopted. The scheme for research sample should not be different for calibration sample. 
    \item Sample can be easily accessed or prepared. Since the main task is to test the system, we should avoid spending much time in preparing the sample. 
    \item Stronger signal amplitude is preferred. This property is dependent on sample phase, spin density, relaxation time and etc. 
\end{itemize}

Nuclei of \ch{^{1}H}, \ch{^{13}C} and \ch{^{15}N} has spin 1/2 and are widely used in NMR experiments. The low matural abundances of \ch{^{13}C} and \ch{^{15}N} (see Tab.\,\ref{tab:isotopesproperity}) imply the difficulty in achieving high spin density of them. In comparison, \ch{^{1}H} is common in compounds. The other advantage of \ch{^{1}H} nucleus, proton, is the high gyromagnetic ratio $\gamma$, which is over 4 times larger than that of \ch{^{13}C} and \ch{^{15}N} nuclei. This brings two benefits to us. The first one is high thermal polarization, which is proportional to $\gamma$ (see Eq.\,\ref{eq:thermalp_approx}). The other benefit is larger frequency band. For the same magnet, the maximal Larmor frequency proton can reach is 4 times larger than that of other nuclei. Hence we are able to search for larger axion mass with protons. 

% \chemfig{^{1}H-^{13}C(-[2]{^{1}H})(-[6]{^{1}H})-[:0]O-[:-72]{^{1}H}}
% \chemfig{^{1}H-^{12}C(-[2]{^{1}H})(-[6]{^{1}H})-[:0]O-[:-72]{^{1}H}}

\begin{table}[htbp]
\centering
\caption{Isotope natural abundance and nuclear gyromagnetic ratio}
\begin{tabular}{ccc}
	\hline
	\multirow{2}{*}{Isotope} & Natural abundance\cite{rumble2019crc} & Nucleus gyromagnetic ratio $\gamma / 2 \pi$ \\
	 & (\SI{}{\%}) & (\SI{}{\MHz. \tesla^{-1}})  \\
	\midrule
	\ch{^{1}H} & \SI{99.9850(70)}{} & \SI{42.577 482 5}{}\cite{rumble2019crc} \\
	\midrule
	\ch{^{3}H} & - & \SI{45.415}{} \\
	\midrule
	\ch{^{3}He} & \SI{0.000137(3)}{} & \SI{32.434 1025}{}\cite{rumble2019crc} \\
	\midrule
	\ch{^{13}C} & \SI{1.07(8)}{} & \SI{10.707 944 6}{} \cite{CAVANAGH20071}\\
	\midrule
	\ch{^{15}N} & \SI{0.368(7)}{} & \SI{4.317 873 6}{} \cite{CAVANAGH20071}\\
	\midrule
	\ch{^{129}Xe} & \SI{26.44(24)}{} & \SI{11.860 156 3}{} \cite{https://doi.org/10.1002/mrc.4191}\\
	\midrule
	\ch{^{131}Xe} & \SI{21.18(3)}{} & \SI{3.515 769 2}{} \cite{https://doi.org/10.1002/mrc.4191} \\
	\bottomrule
\end{tabular}
\label{tab:isotopesproperity}
% \footnotesize{}
\end{table}

Since proton is chosen for the calibration phase of experiment, the next task is to choose the optimal molecule. In Tab.\,\ref{tab:samplesproperity} were included hydrogen (\ch{H2}), ammonia (\ch{NH3}), methanol (\ch{CH3OH}), ethanol (\ch{CH3CH2OH}) and tetramethylsilane (\ch{Si(CH3)4} or TMS) as potential NMR calibration sample. 
Concerning the accessibility, methanol, ethanol and TMS are in liquid phase at room temperature and pressure of atmosphere. While hydrogen and ammonia are not. 
Besides, methanol has highest spin density among them. Therefore we choose it for most of our experiments. Note that we adopt both regular \ch{^{12}C} methanol and \ch{^{13}C}-enriched one for specific purposes. 

%STP - Standard Temperature and Pressure - is defined by IUPAC (International Union of Pure and Applied Chemistry) as air at 0 oC (273.15 K, 32 oF) and 10^5 pascals (1 bar).

%NTP - Normal Temperature and Pressure

%NTP is commonly used  as a standard condition  for testing and documentation of fan capacities:

%NTP - Normal Temperature and Pressure - is defined as air at 20oC (293.15 K, 68oF) and 1 atm (101.325 kN/m2, 101.325 kPa, 14.7 psia, 0 psig, 29.92 in Hg, 407 in H2O, 760 torr). Density 1.204 kg/m3 (0.075 pounds per cubic foot)

%At these conditions, the volume of 1 mol of a gas is 24.0548 liters.
% \begin{table}[ht]
\begin{sidewaystable}[p]
\centering
\caption{Properties of chemicals for NMR\cite{NISTChemistryWebBook}}
\begin{tabular}{ccccccc}
	\hline
	\multirow{2}{*}{Chemical$^{\tiny \cir{1}}$} & Melting point & Boiling point& Density$^{\tiny \cir{2}}$ & Molar mass & Spin density & $\mu / \mu_N$ $^{\tiny \cir{5}}$\\
	 & (\SI{}{\kelvin}) & (\SI{}{\kelvin}) & (\SI{}{\gram. \cm^{-3}}) & (\SI{}{\gram. \mole^{-1}})&(\SI{}{\mmol.\cm^{-3}})& \\
	\midrule
	{\ch{H2} (liquid)} & \multirow{2}{*}{13.96} & \multirow{2}{*}{20.32} & {\SI{7.71e-2}{}} & \multirow{2}{*}{\SI{2.02}{}} & \SI{76.47}{} & \multirow{8}{*}{\SI{+2.79}{}} \\
	{\ch{H2} (gas)} &  &  &{\SI{1.32e-3}{}} &  & \SI{1.31}{}& \\
	\cline{1-6}
	{\ch{NH3} (liquid)} & \multirow{2}{*}{195.49} & \multirow{2}{*}{\SI{239.57}{}} & {\SI{0.734}{}} & \multirow{2}{*}{\SI{17.03}{}} & \SI{86.19}{}&  \\
	{\ch{NH3} (gas)} &  &  & \SI{8.79e-4}{} &  & 0.103 & \\
	
	\cline{1-6}
	\ch{H2O} (liquid) & \SI{273.16}{} & \SI{372.76}{} & \SI{1.00}{} & \SI{18.02}{} &\SI{111.01}{} &  \\
	\cline{1-6}
	\ch{CH3OH} (liquid) & \SI{175.61}{} & \SI{337.30}{} & \SI{0.904}{} & \SI{32.04}{} & \SI{112.87}{} & \\
	\cline{1-6}
	\ch{CH3CH2OH} (liquid) & \SI{159.01}{} & \SI{351.38}{} & \SI{0.789}{} & \SI{46.07}{} & \SI{102.82}{}&  \\
	\cline{1-6}
	\ch{Si(CH3)4} & \SI{172.00}{} & \SI{299.70}{} & \SI{0.648}{} & \SI{88.22}{} & \SI{88.14}{} &  \\
	\midrule
	{\ch{^{129}Xe} (liquid)$^{\tiny \cir{3}}$} & \multirow{3}{*}{\SI{161.40}{}} & \multirow{3}{*}{\SI{164.82}{}} & \multirow{3}{*}{\SI{2.94}{}} & \multirow{3}{*}{\SI{131.29}{}} & \SI{5.92}{} & \multirow{3}{*}{\SI{-0.778}{}}  \\
	{\ch{^{131}Xe} (liquid)$^{\tiny \cir{4}}$}&  &  &  &  & \SI{4.76}{} & \\
	{\ch{^{129}Xe} (\SI{10}{\torr} gas, pure)$^{\tiny \cir{5}}$}&  &  &  &  & \SI{3.88e-3}{} & \\
% 	
%   \midrule
% 	\ch{^{129}Xe}$^a$ & \SI{161.40}{} & \SI{165.051}{} & \SI{2.942}{}$^c$ & \SI{131.2930}{} & \SI{5.92}{} \\
% 	\midrule
% 	\ch{^{131}Xe}$^b$ & \SI{161.40}{} & \SI{165.051}{} & \SI{2.942}{}$^c$ & \SI{131.2930}{} & \SI{4.76}{} \\
	\bottomrule
\end{tabular}

\label{tab:samplesproperity}
\footnotesize{
$^{\tiny \cir{1}}$ pressure at \SI{1}{\bar} if not otherwise noted. 
$^{\tiny \cir{2}}$ at melting point mentioned in the table. 
%$^{\tiny \cir{1}}$ 1.21 MPa, 21 K. $^{\tiny \cir{2}}$ at STP. 
$^{\tiny \cir{3}}$ $^{\tiny \cir{4}}$ in natural-abundance Xenon. 
$^{\tiny \cir{5}}$ Nuclear magnetic moment in units of the nuclear magneton $\mu_N$. }
\end{sidewaystable}
% \end{table}